% Plantilla maestra — Seitz 1903, II. Disputatio Luther-Eck (prueba pp. 55-64)
% Compilar: xelatex plantilla-maestra.tex
\documentclass[11pt,a4paper]{article}
\usepackage{fontspec}
\setmainfont{Libertinus Serif}
\usepackage[lmargin=3cm,rmargin=3cm,top=2.5cm,bottom=2.5cm]{geometry}
\usepackage{setspace}
\setstretch{1.08}
\setlength{\parskip}{0.4em}
\setlength{\parindent}{1.2em}

\title{Der authentische Text der Leipziger Disputation (1519)\\II.\ Disputatio Luther--Eck (Seitz 1903, pp.\ 55--64, prueba)}
\author{Transcripci\'on diplom\'atica de prueba}
\date{}

\begin{document}
\maketitle
\noindent\textit{Fuente: Seitz 1903, test pp.\ 1--10 (= impresas 55--64). Fragmento(s) sin pre\'ambulo, ensamblados por \texttt{\textbackslash input}.}

\bigskip
II.

In nomine domini Amen\footnote{Diese Vorrede findet sich nur in P, vgl.\ oben p.\ 4.}. Anno ab eodem incarnato M.\ D.\ XIX.\ indictione septima, die vero lunae, quarta mensis Julii, pontificatus sanctissimi in Christo patris et domini nostri, domini Leonis divina providentia Papae decimi anno septimo, hora septima antemeridiana, postquam per aliquot dies, nempe a vicesima septima Junii, celebrata esset disputatio inter egregios et praeclaros viros ac dominos Johannem Eccium, canonicum Eystetensem et procancellarium Ingolstatensem, et Andream Rodolphi Bodenstein Carolostadium, canonicum et archidiaconum Wittenbergensem, artium decretorum et theologiae doctores, in materia de libero arbitrio: coepta est disputatio inter magnificos D.\ doctores D.\ Johannem Eccium praefatum et dominum Martinum Lutherum Augustinianum Wittenbergensem in nostra notariorum publicorum et testium infrascriptorum praesentia in hunc qui sequitur modum. Sed quia egregius dominus doctor Martinus in suae disputationis exordio ad priores praefatorum doctorum protestationes se retulit, ideo eaedem ex praecedenti disputatione transsumptae ad verbum sequuntur huius disputationis initio praeinsertae\footnote{Hier folgen in P die Protestationen Karlstadts und Ecks, vgl.\ oben p.\ 14 f.\ [W.\ A.\ 254]}.

Dominus D.\ Martinus in hunc modum orsus est\footnote{U Mc haben hier: Incipit disputatio excellentium theologorum, Johannis Eccii et Martini Lutheri Augustiniani, quae coepta fuit quarta die Julii M.\ D.\ XIX.\ hora septima.\ Protestatio patris M.\ Lutheri [Mc doctoris Martini].}.

'In nomine domini Amen\footnote{U ἀρχήν.}. Protestationem utriusque egregii domini et\footnote{fehlt in M} Andreae Carolostadii et Johannis Eccii amplector et sequor. Hoc unum addo, quod pro reverentia summi pontificis et Romanae ecclesiae libens\footnote{M videns.} hanc materiam non necessariam et mire invidiosam praetermisissem, nisi per propositionem\footnote{M protestationem.} E.\ D.\ Doctoris Johannis Eccii in eam pertractus fuissem.

Doleo etiam eos non adesse, quos maxime oportuit\footnote{M oportet.}, Qui cum et privatim et publice toties me crimine haereseos\footnote{M in haereseos.} prophanarunt, nunc cum instet cognitio causae se subtraxerunt, haereticae pravitatis inquisitores dico, qui fraternam monitionem\footnote{M Lücke, Mc paternam admonitionem.} et doctrinam postposuerunt criminationibus suis.

[255] Dominus doctor Eccius\footnote{U Scopus. De potestate, immo de primatu Romani pontificis (Uv a. R. Additio). Eccii protestatio, Mc [rot] = U.}.

In nomine tuo, dulcis Jesu\footnote{fehlt in M.}.

Antequam in harenam descendam protestor coram, vobis illustrissimis, nobilibus, magnificis et excellentissimis dominis omnia per me dicenda et dicta, inprimis esse subiecta iudicio primae sedis et domini in ea sedentis, deinde\footnote{M dehinc.} quorumlibet aliorum, quorum est errantes corrigere et ad veritatis cognitionem reducere. Et quia reverendus pater in praefatione sua, quasi se excusando, testatur se libenter 'hanc' obmisisse\footnote{U omisisse.} materiam pro summi pontificis reverentia, nisi per meam pertractus fuisset propositionem: at meminerit reverendus pater: nisi prius ipse in resolutorio ante tempora Sylvestri negasset Romanum\footnote{M Romanam.} pontificem aliis\footnote{M aliis esse.} superiorem, non fuisset necessarium me tredecimam illam posuisse propositionem. Et quod in actis coram legato\footnote{M legatis.} sedis apostolicae beatum Pelagium pontificem evangelica scripta\footnote{M dicta.} torsisse causatur. Qui tamen maxime omnium, ad sanctorum patrum intentionem verba Christi accepit. Frustra ergo reverendus pater in me causam torquere nititur, cui ipse non semel occasionem praestitit: sed iam ambagibus resectis\footnote{P reiectis.} principale institutum, deo duce, aggrediamur. Reverende pater, vestra conclusio XIII. meae contradicens, fatetur Romanam ecclesiam' esse' aliis superiorem iuxta frigidissima Romanorum pontificum decreta, intra quadringentos annos\footnote{M quadringenta annis.} nata. Quibus obstare dicitis, textum sacrae scripturae et historias approbatas mille et\footnote{et fehlt in M.} centum 'annorum'. Contra quam sic oppono.

Monarchia et unus principatus in ecclesia dei est de iure divino et a Christo institutus. Quare textus sacrae scripturae vel historiae approbatae ei non adversantur, quoniam ecclesia illa militans, quae est velut unum corpus iuxta divi Pauli sententiam, est instituta et facta ad imaginem ecclesiae triumphantis: in qua est una monarchia
% empalme: la página sigue "omni-" + "bus" = "omnibus" (ver test-p02) -->

omnibus per ordinem dispositis, usque ad unum caput, scilicet deum. Quare et talis ordo a Christo in terris est constitutus, cum fateatur Johannis 5 [v. 19] filium non facere quicquam, nisi quod viderit patrem facientem. Quare non de caelo est, qui capiti subesse recusat. Sicuti non de caelo, sed de lucifero est, qui deo non vult subiici. Quae omnia latissime confirmare possum\footnote{U confirmari possunt.} maxime per sacram illam animam beatum Dionysium Areopagitam, libro de ecclesiastica hierarchia, ubi ait: Nostra enim hierarchia a deo traditis ordinibus sancte disposita, sanctis [256] et caelestibus hierarchiis conformis est. Ita Gregorius Nazianzenus in Apologetico, sacrosancta mysteria, ait ad\footnote{fehlt in M.} instar caelestis imaginis\footnote{P maginis.} celebrari. Per quae\footnote{M terque [gestrichen], Mc per quae.} utique in terris caelestibus ordinibus associamur. Nam quod monstrum esset ecclesiam esse acephalam\footnote{M acevolam, U a. R. Id est, sine capite.}: quod omnes ferme haeretici moliti sunt, ut sanctus Cyprianus ad Rogatianum et Puppianum\footnote{M Rogatianum et Cupianum, U rogationum et Puppianum.} innuit, ut sic debilitato capite errores suos et virus impune possent mentibus hominum inculcare. Et haec fuit ratio principalis cum aliis annexis, propter quam Parisiense felix\footnote{M Parisiensis velut.} studium Johannem Tornacensem damnavit, negantem primatum ecclesiae Romanae. Sic et error fuit Vuiclefficus\footnote{M wiewerbiticus, Mc wicleffiticus.} Romanam ecclesiam aliis non esse superiorem de lege evangelica\footnote{U evangelii.}.

D. doctor Martinus.

Quando dominus doctor arguit esse omnino\footnote{M U omnino esse.} caput ecclesiae universale, optime facit. Et si est aliquis, qui privato\footnote{M a. R. praefato.} pacto convenerit cum\footnote{fehlt in M.} dominatione sua ad defendendum contrarium: surgat in medium, ad me nihil pertinet.

D. doctor Eccius.

Quia reverendus pater dicit nihil ad se pertinere de contrario illius, quod intendebam probare, de iure divino esse monarchiam in ecclesia militante: sicut et triumphante, in quo ego\footnote{fehlt in U.} eum collaudo, quod et in hoc divo Johanni in Apocalypsi consentit: Vidi civitatem sanctam novam, descendentem etc. At\footnote{M ad, Mc at.} propius ad rem accedentes: Si ecclesia militans non fuit sine monarchia\footnote{U monarcha.}, vellem audire, quis esset iste monarcha\footnote{M ista monarchia.} alius aut unquam fuisset, nisi Romanus pontifex, aut quae\footnote{M atque (statt ant quae).} alia prima sedes, nisi sedes Petri et eius successorum, iuxta illud beati Cypriani, qui in epistola tertia\footnote{M P U secunda.} ad Cornelium Romanum pontificem scribens, contra Novatianos subdole Romam accedentes, inquit: Post ista adhuc insuper, pseudoepiscopo\footnote{P U cedo episcopo, M sedo episcopo.} sibi ab haereticis constituto\footnote{W. A. setzt fälschlich hinzu Romam und Anm. „Romam fehlt"; corrigiert am Ende des 2. Bandes.} navigare audent et ad Petri cathedram\footnote{U. a. R. non ad Petris cabellum.} atque ecclesiam principalem, unde sacerdotalis unitas exorta est, a schismaticis et prophanis literas ferre, nec\footnote{U ne.} cogitare eos esse Romanos, quorum fides ab Apostolo laudata est, ad quos perfidia habere non possit accessum. Et Hieronymus idem testatur contra Luciferianos, ecclesiae salus, ait, in summi sacerdotis dignitate pendet: cui si non exors\footnote{M exsors.} quaedam et ab omnibus eminens detur potestas, tot in ecclesiis efficientur\footnote{M tot in ecclesiis tot efficienter (so!)} schismata, quot sacerdotes. Et quod ille summus sacerdos sit Romanus, liquet apud eundem beatum Hieronymum in duabus epistolis ad Damasum papam, cuius ferme singula verba ad institutum faciunt\footnote{M a. R. facio.}. Verum brevitatis studio illa signemus: Cum successore piscatoris et discipulo\footnote{U discipuli.} Christi loquor. Ego nullum praemium sequens nisi Christum, beatitudini tuae id est cathedrae [257] Petri consocior\footnote{M consocio.}, super illam petram fundatam ecclesiam scio; et infra: Quicumque tecum non colligit ille dispergit. Ex quibus quisque bonus christianus facile coniicit et intelligit\footnote{et intelligit fehlt in U.} unitatem sacerdotalem a Romano fluere pontifice, et illam fuisse semper cathedram principalem et omnibus aliis praelatam\footnote{M praelatis.}, et quod illa sit petra, sicut Hieronymus dicit se\footnote{M se dicit.} scire, super quam fundata sit ecclesia. Vel reverendus pater alium monarcham\footnote{M aliam monarchiam.} ecclesiae priscis temporibus assignet\footnote{M U assignet (M assignat) alium monarcham etc.}.

Dominus D. Martinus.

Monarchiam ecclesiae militantis prorsus confiteor eiusque caput non hominem sed Christum ipsum idque auctoritate divina prima ad Corinthios 15 [v. 25]: Oportet autem\footnote{fehlt in M U.} illum regnare, donec ponat omnes inimicos sub pedibus eius. Et paulo ante: Deinde finis, cum tradiderit regnum deo et patri et evacuaverit omnem principatum. Quod beatus\footnote{fehlt in U.} Augustinus libro I.\footnote{libro I fehlt in M.} de trinitate cap. ultimo\footnote{fehlt in M, ultimo cap U.} exponit de regno Christi praesentis temporis. Ita videlicet, quod Christus caput ecclesiae per fidem transferet nos, qui regnum eius sumus, ad\footnote{ad P M Uv, per U.} speciem. Sic Matth. ultimo: Ecce ego vobiscum sum usque ad consummationem saeculi. Item actuum 9 [v. 4]: Audivit Paulus de caelo: Saule, Saule,

quid me persequeris? ubi iterum Augustinus dicit, caput loqui pro membris suis\footnote{fehlt in M.}. Quare prorsus audiendi non sunt, qui Christum extra ecclesiam militantem trudunt in triumphantem\footnote{Mc ecclesiam triumphantem.}, cum sit regnum fidei\footnote{M filii.}. Hoc est: quod caput nostrum non videmus et tamen\footnote{M attamen.} habemus. Juxta illud Psalmi 122\footnote{Mc obsequendam statt ad obsequendi.}: Illic sederunt sedes in iudicio super domum David. Multae scilicet\footnote{Uv sunt.} sedes, in quibus unus sedet Christus\footnote{Mc scilicet Christus.}. Sedes videmus, non sessorem vel regem. Deinde\footnote{M U proinde.} ad auctoritates domini\footnote{M U egregii domini.} D. veniendo\footnote{M veniendo ad auctoritates etc.} quando asserit in ecclesia militante\footnote{M triumphante, Mc militante.} esse unum principatum divino iure et a Christo institutum\footnote{M instituto.}, sua dicit, sed nihil probat. Nam prima eius auctoritas Pauli, nimirum ad\footnote{fehlt in P U.} Ephes. 4 [v. 15], quae Christum caput ecclesiae dicit, pro me, contra eum facit\footnote{U valet.}, cum certissime loquatur ibi\footnote{M ibi loquatur.} de ecclesia\footnote{M ecclesiae.} militante et Christum appellet\footnote{M U appellat.} eius caput. Est etiam\footnote{M Est est.} eadem sententia, huic contraria\footnote{M U contraria huic.} I. ad Corinth. 3 [v. 5; 1 v. 13]: Quid est Apollo? quid Cephe? quid Paulus? divisus est Christus?\footnote{U add. etc.} Ubi manifeste prohibet aliud caput quam Christum.

Secunda eius\footnote{fehlt in M.} auctoritas Johannis 5.: Non potest\footnote{M poterit.} filius facere quicquam\footnote{M quicquam facere.}, nisi quod\footnote{fehlt in M.} viderit patrem facientem, neque de ecclesia militante neque triumphante loquitur, sed omnium doctorum sententia de aequalitate sua cum patre, scilicet\footnote{U sed, Uv scilicet.} quod nihil possit facere pater, quin idem possit filius\footnote{M quod nihil faciat nec facere possit, quod non idem facit filius, U quod nihil facit aut possit facere pater, quin idem possit et filius.}. Omitto\footnote{M 'Omitto' omitto.} quod dixit\footnote{fehlt in M.} non esse de caelo, qui capiti subesse recuset\footnote{U recusat.}, et esse de lucifero, 'qui\footnote{M est, qui.} deo nolit\footnote{U non vult.} subiici'; quia sicut praecedentia sunt male\footnote{U male sunt.} assumpta, ita\footnote{M sic, Mc ita.} haec\footnote{U et haec.} male illata.

Tertia, 'quae et Dionysii', nihil contra nos. Non enim negamus hierarchiam ecclesiasticam\footnote{M coelestem, Mc ecclesiae.}, sed de capite disputamus monarchiae, non hierarchiae\footnote{M monarchia, non hierarchia.}.

[258] Quarta\footnote{U Quarto.}, quae est\footnote{fehlt in M U.} Gregorii Nazianzeni inducitur\footnote{M dicitur Mc additur (adducitur U), al. inducitur.}, quod per sacrosancta mysteria sociamur\footnote{M sociamur.} ordinibus caelestibus\footnote{Mc coelestibus ordinibus, U sanctis ordinibus coelestibus.} etc., 'a quovis grammatico intelligitur' nihil sonare neque de monarchia neque 'de capite'\footnote{M capite hierarchia, letztes Wort gestrichen.}. Jam quod addidit\footnote{M nam quod adducit} esse monstrum, si\footnote{M quod, Mc si.} ecclesia acephala\footnote{M A scevola.} esset, fateor\footnote{M (Mc) 'fateor' fateor.}; sed hoc caput nec ipse\footnote{fehlt in M.} D. doctor aliud dare\footnote{M dari.} possit quam Christum, quod probo evidenter: quia si suum caput, quod Romanum pontificem appellant\footnote{M vocat romanum pontificem, U Romanum pontificem appellat.}, moritur, ut\footnote{Mc quia.} est homo, iam ecclesia est acephala\footnote{M acephola.}. Si autem Christus interim caput est ecclesiae, donec alius eligatur, non minus monstrum\footnote{M U est.}, Christum 'cedere' vivo pontifici et succedere mortuo.

Quinta beati Cypriani, qui haereticos insectatus\footnote{M insectatur, Mc insectatus.}, quod debilitato capite amoliti\footnote{U moliti.} sunt errores suos impune hominibus inculcare: prorsus nihil facit ad propositum. Loquitur\footnote{M U de quolibet.} enim non de Romano capite, sed quolibet\footnote{M U de quolibet.} capite cuiuslibet\footnote{U cuiusque.} 'episcopatus'. Et\footnote{M Additur et, Mc a. R. Atamen (so!) et.} si egregius\footnote{fehlt in M.} D. doctor voluerit stare auctoritati 'Cypriani', hac hora finiemus disputationem\footnote{fehlt in U.}. Ipse enim Cornelium Romanum\footnote{U a. R. Eccius dixit non fiet.} pontificem nunquam aliter salutavit quam\footnote{U salutat nisi.} charissimum fratrem. Deinde describens episcoporum electiones et confirmationes\footnote{M electionem et confirmationem.} per multas epistolas ex divinis literis efficacissime\footnote{M longissime, Mc efficacissime.} probat, ad plebem pertinere et vicinos duos episcopos\footnote{M suos episcopos duos.} vel tres, quemadmodum et sancitum est in sacratissimo concilio Niceno\footnote{U Niceno concilio; M add. ut beatus Augustinus allegat, gestrichen, weil vom Disputator widerrufen.}. Quin idem beatus martyr, ut beatus Augustus allegat libro 2 ca. 2 de baptismo\footnote{Mc add. parvulorum.}, sic dicit\footnote{M dicens.}: Neque enim quisquam nostrum episcopum se essee piscoporum constituit\footnote{M constituat.}, aut tyrannico terrore\footnote{U errore.}, ad 'obsequendi\footnote{Mc obsequendam statt ad obsequendi.} necessitatem' collegas suos adigit\footnote{M addidit.}, quando habet omnis episcopus pro licentia libertatis et potestatis suae arbitrium proprium, tanquam ab alio iudicari non possit\footnote{M posse.}, quomodo nec 'ipse potest alterum iudicare'. Sed expectemus universi\footnote{M omnes.} iudicium domini nostri Jesu Christi\footnote{iudicium domini nostri Jesu Christi fehlt in M, dafür eine Zeile frei.}.

Quod autem dicit de Romana et Petri cathedra ortam esse\footnote{fehlt in M.} sacerdotalem unitatem\footnote{M dignitatem.}, libentissime admitto, quantum pertinet ad ecclesiam occidentalem. Nam revera etiam Romana ecclesia orta est ex Hierosolymitana: haec est proprie matrix\footnote{U matrix proprie, M nur matrix.} omnium ecclesiarum.

Sed nec valet consequentia: Ex Romana ecclesia orta est sacerdotalis unitas; ergo ipsa est caput et domina omnium\footnote{U add prima.}. Alioquin insuperabiliter concluderet Hierosolymitanam esse caput et dominam\footnote{et dominam fehlt in M.} omnium.

Ultima auctoritas divi Hieronymi non recte inducitur ab egregio domino doctore, etiam si per omnia esset\footnote{M esse, Mc esset.} vera auctoritas Hieronymi\footnote{M beati Hieronymi.}, quia d. doctor intendit probare ecclesiae Romanae potestatem monarchicam iure divino et\footnote{fehlt in U.} a Christo institutam. Hoc verba Hieronymi non habent, quia dicit: Cui\footnote{fehlt in M} si non exors quaedam et\footnote{fehlt in U.} ab omnibus eminens\footnote{M eminus, Mc eminuens.} detur potestas, tot exorientur\footnote{M efficenter (so!), Mc efficerentur, U in ecclesiis efficientur.} in ecclesia schismata quot sacerdotes\footnote{M add sunt.}. Detur (inquit)\footnote{M add non habet.}, hoc est iure humano possit\footnote{M U posset.} fieri, consentientibus ceteris omnibus fidelibus. Nam nec ego hoc nego\footnote{M et ego non nego.}, si consenserint totius orbis fideles\footnote{M omnes fideles, a. R. omnes totius orbis.} in Romanum\footnote{M Romanam.} vel Parisiensem\footnote{U pariensem.} vel Magdeburgensem vel quemcumque [259], ut esset primus pontifex 'et summus': hunc propter reverentiam totius ecclesiae fidelium sic consentientis\footnote{M consentientes.} habendum esse summum monarcham\footnote{M unam monarchiam, a. R. summum monarchum (so!)}. Hoc autem nec factum est unquam neque fit neque fiet\footnote{M factum est nunquam neque fit neque fiet, a. R. = P U.}, cum ad nostra usque\footnote{M U usque ad nostra.} tempora Graeca ecclesia non consenserit nec tamen habita sit\footnote{M U sit habita.} haeretica. Et hanc esse\footnote{M video esse.} divi Hieronymi sententiam\footnote{U sententiam divi [fehlt in M] Hieronymi.} probo ex epistola ad Euagrium\footnote{M ephasum.}, ubi dicit: Ubicumque episcopus fuerit sive Romae sive Eugubii seu Constantinopoli sive Regii sive Alexandriae sive Tanis, eiusdem meriti et eiusdem sacerdotii est. Potentia divitiarum et humilitas paupertatis vel sublimiorem vel inferiorem facit. Ceterum omnes apostolorum successores sunt\footnote{U a. R. Martinus: Das sie den Hiero: nicht lang verbrant haben.}. Recitatur haec epistola in decretis non frigidis\footnote{M vidi.} dist. XCIII. c. legimus. Idem super epistolam ad Titum\footnote{M deditum, Mc Titum.}: Idem\footnote{M Sicut ergo est vor Idem, gestrichen.} est ergo\footnote{M ergo est.} presbyter qui episcopus. Et antequam diaboli instinctu studia in religione fierent et diceretur in populis\footnote{M populo.}: Ego sum Pauli, ego 'Apollo, ego Cephae'\footnote{U Ego sum Pauli, ego sum Cephae.}, communi presbyterorum consilio ecclesiae gubernabantur\footnote{M ecclesia gubernabatur.}. Postquam vero unusquisque 'eos', quos baptizabat, suos esse putabat, in toto orbe decretum est, ut unus de presbyteris electus supraponeretur\footnote{M unus vir super presbyteros superponeretur, U = P, nur superponeretur.}. Et inductis auctoritatibus scripturae dicit in fine: Sicut ergo presbyteri sciunt se ex ecclesiae consuetudine ei, qui sibi praepositus fuerit, esse subiectos, ita episcopi noverint se magis consuetudine quam dispensationis divinae\footnote{U dispositionis dominicae, W. A. Anm. 2 vermutet die richtige L. A.} veritate presbyteris\footnote{M presbyteres.} esse maiores. Ideo quod dominus doctor addidit, a Hieronymo intellectum esse 'summum' pontificem Romanum, dum dicit: Cum successore piscatoris et discipulo\footnote{U discipuli.} Christi loquor et beatitudini tuae id est cathedrae tuae\footnote{U cathedrae Petri, M capite Petro, a. R. catedrae.} consocior\footnote{fehlt in M, U consotius.}, super\footnote{M et super.} illam petram fundatam ecclesiam scio: nihil facit. Non sequitur, consocior huic ecclesiae, ergo haec est sola prima. Nec sequitur, est fundata supra\footnote{M super.} petram, ergo sola est fundata.

Accedit ad hoc decretum concilii Africani dist. XCIX. c. Primae: Primae, inquit, sedis episcopus non appelletur princeps sacerdotum aut summus sacerdos aut aliquid huiusmodi, sed tantum primae sedis episcopus. Universalis autem pontifex nec Romanus appelletur. Quare si divino iure staret monarchia Romani pontificis\footnote{M Romana ierarchia summorum pontificum.}, haec omnia essent haeretica: quod est temerarium asserere. Et in fine audiamus ipsum dominum, qui Luc. 22 [v. 24-26] dicit: Facta est autem contentio\footnote{U contentio discipulorum.}, quis eorum videretur esse maior. Dixit autem eis: Reges gentium dominantur eorum\footnote{U earum.}, et qui potestatem habent super eos, benefici\footnote{M heretici.} vocantur: vos autem non sic, sed qui maior est vestrum, fiat sicut minor\footnote{M add. Ecce hoc (haec?) est mea sententia.}.

D. D. Eccius\footnote{In M beginnt mit diesem Abschnitt eine neue Seite Bl. E 3a. Auf der vorangehenden findet sich ein andrer (von Eck widerrufener?) Anfang: „Reverendus pater debens;" dann folgt ein leerer Raum von etwa 4 Zeilen.}.

Reverendus D. pater satis instructus descendit in harenam, materiam suam habens bono ordine contextam in libello impresso et per eum composito [260]. Quare illustres\footnote{U illustrissimae.} dominationes vestrae, magnificentiae et dominationes\footnote{M et magnificentiae statt magnificentiae et dominationes.} veniam dabunt Eccio, multo iam tempore aliis negotiis occupato, si non tam rotunde et accurate e vestigio\footnote{PU evestigio} tanta cumulare poterit, quanta\footnote{P quantam.} reverendus pater iam congessit. Venit enim disputaturus, non librum editurus. Verum quid reverendus pater dixerit, per ordinem dispiciamus.

Primo omnium caput esse ecclesiae Christum probare intendit: quod tamen fuerat supervacaneum, cum nullus hoc negare praesumat, nisi qui Antichristus sit\footnote{fehlt in M.}. Demiror\footnote{MU miror.} tamen vehementer, quod hoc

non perpendit, sicut pollicetur in actis coram legato sedis apostolicae posse se\footnote{M se posse.} praestare quendam iuristotheologum\footnote{M iuristam theologum, U iuristatheologum.}, plura posse esse\footnote{U esse posse, M om. posse.} capita subordinata. In quo mystici capitis ratio aut etiam symbolici\footnote{M et symbolici, U aut symbolici.} a ratione capitis naturalis deficit. Quare praeter caput Christum et aliud oportere quaeri caput in ecclesia statim probabitur. Neque ei suffragatur, quod apostolum inducit I. ad Cor. 3 [v 5; 1 v. 13]: Divisus est Christus etc.\footnote{fehlt in U.} Nam etsi Paulus ibi Petri meminerit, non tamen falsum dixit beatus Hieronymus libro 1. contra Jovinianum, mihi columna 18: Unus eligitur (loquitur de Petro) ut capite constituto schismatis\footnote{M ut schismatis.} tollatur occasio. Clare appellat Petrum caput in ecclesia constitutum\footnote{U esse constitutum.}. Sed haec missa faciamus\footnote{M U facimus.}, solutiones ad nostra inducta repulsuri\footnote{M refusuri.}.

Primo cum respondit\footnote{M U respondet.} ad illud Johannis 5: Neque enim potest filius\footnote{M U filius potest.} facere quicquam, nisi quae viderit patrem facientem, dixit secundum sanctos patres exprimi hic\footnote{fehlt in M.} aequalitatem patris et filii. Verum legat attentius quaeso reverendus pater beatum et inadulabilem\footnote{M inadulabilem, Mc adulabilem.} patrem Bernhardum libro 3. de consideratione ad Eugenium, ubi de forma ecclesiae loquens et eam esse de iure divino probans, nostram formavit rationem\footnote{M formavit probationem, U rationem formavit.}, mihi columna 7.\footnote{mihi col. VII fehlt in M.} Nec vilem reputes\footnote{U reputans, Uv reputes.} formam hanc, quia in terra est: exemplar habet in coelo. Neque enim filius potest facere quicquam, nisi quae viderit patrem facientem. Praesertim cum ei sub Moysi nomine dictum sit: Vide, omnia\footnote{M ne omnia.} facias secundum exemplar, quod tibi in monte monstratum est. Viderat hoc, qui dicebat\footnote{U dicebar.}: Vidi civitatem sanctam\footnote{U add. etc.}. Ego enim propter similitudinem dictum reor: quod sicut illic Seraphim et Cherubim\footnote{M cherubim et seraphim.} et 'ceteri' quique\footnote{U ceteri qui; M ceteri, constitutique, gestrichen constituti.} usque ad angelos et archangelos ordinantur sub uno capite deo, ita hic quoque sub uno summo pontifice primates vel patriarchae, episcopi, presbyteri\footnote{M Archipresbyteri statt episcopi, presbyteri; U Archiepiscopi, episcopi, presbyteri.} vel abbates et reliqui in hunc modum. Tunc\footnote{fehlt in M.} 'subdit Bernhardus': Non est parvi pendendum, quod et deum habet\footnote{U habeat.} auctorem, et de caelo ducit originem. Quis iam non intelligit\footnote{U intelligat.}, hanc ecclesiasticam hierarchiam secundum Bernhardum a Christo institutam? Et velut Christus\footnote{M U deus.} caput est in caelo, ita summus pontifex caput est\footnote{M U sit.} in militante ecclesia, neutiquam\footnote{M nec dicam.} Christum excludendo, cuius se fatetur vicarium. Sed ad ratiunculam plebeiam\footnote{U pleiam, Uv plebeiam.}, quam inducit\footnote{U induxit.}, ecclesiam manere acephalam\footnote{M ad Scepolam.}, [261] papa mortuo, nisi ridicule diceremus Christum cedere papae vivo et succedere papae mortuo, 'ridicula prorsus ratiuncula', quae\footnote{M et quae.} in re tam seria, inter tot praestantes viros, vix est digna renarrari\footnote{M U enarrari.}, cum ab initio dixerim, caput illud symbolicum a veri capitis 'naturalis'\footnote{M naturalis 2 mal.} ratione in aliquibus convenientiis deficere. Neque Christus, cuius regnum manet in aeternum et sempiternum est eius\footnote{U et cuius sempiternum est.} sacerdotium, cedit papae vel succedit, cum ei sit data omnis potestas in caelo et in terra (Matthaei ultimo) et defuncto papa modo coetus cardinalium, sicut defuncto episcopo capitulum, iura illa teneant\footnote{M U teneat, P Uv teneant.}, usquequo novus pontifex deligatur\footnote{M delegatur.}.

Quod vero secundo loco reverendus pater dicit, Cyprianum loqui de quolibet episcopo, non de Romano pontifice, plurimum miror, cum intelligentia dictorum ex causis dicendi\footnote{M dicendis.} sit accipienda et Cyprianus in locis per me\footnote{M a me.} adductis eos increpet, qui a Cornelio defecerunt, qui certe Romanus fuit pontifex. Ob id sciat reverendus pater me non nudis\footnote{M nullis, Mc nudis.} verbis satiari\footnote{M satiare, Mc satiari.}, ut sophistas pascere solemus. Quae ex Cypriano pro se affert, opponendo arbitror cumulatius efficiet. Nam quod Cyprianus Cornelium fratrem appellat et apostolos fratres fuisse, nemo nescit. Attamen Petrus sicut et eius successor Cornelius caput fuit apostolorum, apex et vertex iuxta beati Dionysii ca. 3. vel. 7. de divinis nominibus sententiam\footnote{fehlt in M (Lücke).}. De electione et a\footnote{fehlt in U (ist in Uv zugefügt).} concilio Niceno et a Cypriano conscripta nec promovent nec impediunt negotium. Et multo minus Augustinus lib. 2. cap. 2. de baptismo parvulorum debebat adduci post Cyprianum. Nam Augustinus arrogantiam et temeritatem eorum, qui per ambitionem et superbiam ad ecclesiasticas praelaturas se ingerunt, reprehendit: quod non debeant se constituere nec alios cogere ad se constituendum, cum quisque praelatus expectare debeat, ut vocetur, sicut Aaron.

Tertio aliam Cypriani sententiam, ex tertia epistola\footnote{P M U ex secunda epistola.} ad Cornelium, diluturus\footnote{M dilaturus.} respondit unitatem quidem sacerdotalem in occidentali ecclesia a Romana exortam, non in orientali. Obticuit hic reverendus pater, quod Cyprianus praecedenter Romanam ecclesiam appellavit cathedram Petri\footnote{M U petri cathedram.} atque ecclesiam principalem. Sed quid sua solutio possit, medullam verborum inspicienti est manifestarium. Nam reverendus pater more grammatico ortum unitatis sacerdotalis intellexit\footnote{M U intelligit.} quo ad auspicia\footnote{U auspicium.} et initium\footnote{M initia.}, cum profecto Cyprianus voluerit explicare ortum commissionis,

% Siguientes fragmentos: \input{test-p06.tex} etc., unidos sin \newpage salvo corte real de secci\'on.
\end{document}
